% ============================================================
%  01_gioi_thieu.tex  --  Gioi thieu
% ============================================================

\section{Giới thiệu}

\subsection{Bối cảnh bài toán}

Ngôn ngữ C vẫn được sử dụng rộng rãi trong lập trình hệ thống, nhúng và các
ứng dụng đòi hỏi hiệu năng cao. Tuy nhiên, C không cung cấp cơ chế kiểm soát
bộ nhớ tự động, khiến lập trình viên phải tự quản lý toàn bộ vòng đời của vùng
nhớ -- từ cấp phát đến giải phóng. Khi thao tác này thiếu cẩn thận, các lỗi bộ
nhớ nghiêm trọng có thể xảy ra: ghi vượt ranh giới bộ đệm làm hỏng dữ liệu lân
cận hoặc cho phép kẻ tấn công điều khiển luồng thực thi; tiếp tục sử dụng con
trỏ sau khi vùng nhớ đã được giải phóng tạo ra hành vi không xác định có thể bị
khai thác; không giải phóng bộ nhớ sau khi dùng xong dẫn đến cạn kiệt tài
nguyên trong các hệ thống hoạt động lâu dài.

Bài thực hành này được thiết kế trong bối cảnh hệ thống SecureSys --
một hệ thống quản lý truy cập và giao dịch an toàn. Module được chọn là
Inventory -- quản lý kho sách thư viện -- vì đây là module đơn giản,
đủ để minh họa đầy đủ ba lớp lỗi mà không làm phức tạp hoá bài toán. Mỗi cuốn
sách được lưu trong một struct \textit{Book} gồm tiêu đề, tác giả, năm xuất bản
và số lượng bản sao hiện có. Hoạt động cốt lõi -- ghi tiêu đề sách -- là điểm
chèn của cả ba lỗi bộ nhớ.

\subsection{Mục tiêu học tập}

Bài thực hành hướng đến các mục tiêu cụ thể sau:

\begin{itemize}
    \item Phát hiện lỗi bằng phân tích động: Sử dụng AddressSanitizer
          (ASan) -- công cụ được tích hợp trực tiếp trong GCC/Clang -- để phát
          hiện lỗi bộ nhớ tại thời điểm chạy, đọc và phân tích thông báo lỗi
          cùng stack trace do ASan xuất ra.
    \item Nhận diện và phân loại lỗi: Ánh xạ từng lỗi tìm thấy sang
          định danh CWE tương ứng (CWE-121, CWE-416, CWE-401) và giải thích cơ
          chế khai thác tiềm năng.
    \item Sửa lỗi theo nguyên tắc lập trình C an toàn: Thay thế các
          API không an toàn bằng API giới hạn kích thước, kiểm soát quyền sở hữu
          con trỏ và đảm bảo vòng đời bộ nhớ đúng đắn.
    \item Hiểu giới hạn của công cụ: Thực nghiệm và giải thích tại sao
          tràn bộ đệm bên trong cùng một struct không luôn bị ASan phát hiện,
          trong khi tràn trên heap buffer riêng biệt luôn bị bắt.
    \item Đánh giá rủi ro định lượng: Áp dụng mô hình CVSS-lite để
          tính điểm rủi ro, xếp hạng mức độ nghiêm trọng và đề xuất thứ tự ưu
          tiên khắc phục cho năm lỗ hổng.
\end{itemize}

\subsection{Phạm vi bài thực hành}

Bài thực hành bao gồm các thành phần sau:

\begin{itemize}
    \item Mã nguồn C: Hai file song song --
          \textit{c/vulnerable\_inventory.c} chứa ba lỗi cố ý và
          \textit{c/secure\_inventory.c} là bản đã sửa. Cả hai hỗ trợ ba chế độ
          dòng lệnh: \textit{overflow}, \textit{uaf}, \textit{leak}.
    \item Dataset lỗ hổng: File \textit{dataset/vulnerabilities.json}
          mô tả năm lỗ hổng (VULN-001 đến VULN-005) với các thuộc tính CVSS-lite
          (attack vector, complexity, privileges, user interaction, C/I/A impact).
    \item Script đánh giá rủi ro: \textit{python/risk\_assessor.py} đọc
          dataset, tính toán điểm số và xếp hạng mức độ nghiêm trọng.
    \item Bộ kiểm thử tự động: \textit{tests/test\_memory\_safety.py}
          với 8 bài kiểm tra tích hợp cả phát hiện lỗi runtime và kiểm tra tính
          nhất quán của mô hình rủi ro.
\end{itemize}

\subsection{Cấu trúc báo cáo}

Báo cáo được tổ chức theo trình tự thực hành: Phần~\ref{sec:kienthuc} trình bày
kiến thức nền tảng về các lớp lỗi và công cụ. Phần~\ref{sec:dacta} đặc tả hệ
thống và luồng dữ liệu. Phần~\ref{sec:hopden} trình bày kết quả phát hiện lỗi
bằng ASan kèm stack trace. Phần~\ref{sec:hoptrang} trình bày bản vá code và
đánh giá rủi ro. Phần~\ref{sec:tudong} tổng hợp kết quả kiểm thử tự động.
Phần~\ref{sec:ketluan} thảo luận và kết luận.
